\subsection{সমন্বিত CQ: হিসাব, কারণ ও নির্বাচন}
এই অনুশীলনে বিভিন্ন ধারণা একসঙ্গে প্রয়োগ করতে হবে: হিসাব, কারণ ব্যাখ্যা এবং পরিস্থিতির উপযোগী প্রযুক্তি নির্বাচন।

\subsubsection{CQ–১৩: Character ও Block-এর দক্ষতা}
একটি যন্ত্র ১০০টি ক্যারেক্টার পাঠায়। প্রতিটিতে ৮টি ডেটা বিট আছে। পদ্ধতি A-তে প্রতিটি ক্যারেক্টারের জন্য ১টি স্টার্ট ও ২টি স্টপ বিট ব্যবহৃত হয়; প্যারিটি নেই। পদ্ধতি B-তে পুরো ব্লকের হেডার ও ট্রেলার মিলিয়ে ৪ Byte বাড়তি তথ্য আছে। উভয় পদ্ধতিতে লাইন হার ১১০০ bps।
\begin{compactenum}
\item ক. বিট সিনক্রোনাইজেশন কী?\quad খ. স্টার্ট ও স্টপ বিট প্রয়োজন কেন?
\item গ. A-এর দক্ষতা এবং উপকারী ডেটা হার নির্ণয় করো।
\item ঘ. B-এর দক্ষতা গণনা করে সিদ্ধান্ত দাও।
\end{compactenum}

\subsubsection{CQ–১৪: Office-এর Wireless নির্বাচন}
প্রতিষ্ঠানের ফোন LTE ব্যবহার করে। অফিসের ল্যাপটপ 802.11 সংযোগে যুক্ত। পাশের বড় এলাকায় 802.16 সংযোগ পাওয়া যায়। ৪০ km দূরের শাখার সঙ্গে যোগাযোগ প্রয়োজন। একটি ছোট স্বয়ংক্রিয় যন্ত্র অনিয়মিত বিরতিতে ক্যারেক্টার পাঠায়।
\begin{compactenum}
\item ক. রোমিং কী?\quad খ. অনিয়মিত বিরতিতে পাঠানো ক্যারেক্টারে ফ্রেমের সীমানা নির্ধারণ প্রয়োজন কেন?
\item গ. ফোন ও স্থানীয় নেটওয়ার্কের প্রযুক্তি এবং ব্যবহার ব্যাখ্যা করো।
\item ঘ. শাখার সংযোগের জন্য 802.11 ও 802.16-এর মধ্যে উপযোগী প্রযুক্তি নির্বাচন করো।
\end{compactenum}

\subsubsection{CQ–১৫: Cable ও Network Device}
হাসপাতালের কক্ষগুলোর কম্পিউটার চার জোড়া তামার তারবিশিষ্ট কেবল দিয়ে সুইচে যুক্ত। ২ km দূরের ভবনে EMI-এর প্রভাবমুক্ত সংযোগ প্রয়োজন। ইন্টারনেট ব্যবহারে রাউটার এবং ভিন্ন প্রোটোকলের পুরোনো যন্ত্র যুক্ত করতে অন্য একটি ব্যবস্থা ব্যবহার করা হবে।
\begin{compactenum}
\item ক. NIC কী?\quad খ. সুইচ কি সব ফ্রেম শুধু একটি পোর্টে পাঠায়?
\item গ. দুটি মাধ্যম শনাক্ত করে ব্যবহারের কারণ লেখো।
\item ঘ. রাউটার, প্রোটোকল গেটওয়ে ও রিপিটারের ভূমিকা তুলনা করো।
\end{compactenum}

\subsubsection{CQ–১৬: Topology, Port ও Failure}
ল্যাবের ৬টি কম্পিউটার কেন্দ্রীয় সুইচে যুক্ত। ছয়টি জরুরি সার্ভারের প্রতিটি জোড়ায় পৃথক দ্বিমুখী লিংক রাখা হবে। ক্যাম্পাসের বিভাগীয় সুইচগুলো স্তরভিত্তিকভাবে কেন্দ্রের সঙ্গে সংযুক্ত।
\begin{compactenum}
\item ক. ভৌত টপোলজি কী?\quad খ. স্টারে একটি কেবল নষ্ট হওয়া ও কেন্দ্রীয় সুইচ বিকল হওয়ার প্রভাব আলাদা কেন?
\item গ. সার্ভারগুলোর জন্য প্রয়োজনীয় লিংক ও পোর্টের সংখ্যা গণনা করে চিত্র আঁকো।
\item ঘ. ল্যাব ও ক্যাম্পাসে Mesh এবং Star/Tree-এর উপযোগিতা মূল্যায়ন করো।
\end{compactenum}

\subsubsection{CQ–১৭: Cloud Model ও হারানো File}
কলেজ প্রস্তুত অনলাইন ডকুমেন্ট সফটওয়্যার ব্যবহার করে। নিজস্ব অ্যাপ্লিকেশন একটি পরিচালিত প্ল্যাটফর্মে চলে। ফলাফলের Private Cloud অনুমোদিত ডেটা Public Admission Cloud-এর সঙ্গে সমন্বিতভাবে ব্যবহার করে। একজন শিক্ষার্থীর ল্যাপটপে ফাইল মুছে ফেলার পর ক্লাউডের সিঙ্ক্রোনাইজড কপিও মুছে গেছে।
\begin{compactenum}
\item ক. Measured Service কী?\quad খ. সিঙ্ক্রোনাইজেশনকে ব্যাকআপ বলা যথেষ্ট নয় কেন?
\item গ. সেবা মডেল ও স্থাপন মডেল শনাক্ত করো।
\item ঘ. গোপনীয়তা, খরচ এবং ফাইল পুনরুদ্ধারের পরিকল্পনা দাও।
\end{compactenum}
